% Generated draft: requires mathematical and source review.
\documentclass[11pt]{article}
\usepackage[margin=25mm]{geometry}
\usepackage{fontspec}
\setmainfont{Latin Modern Roman}
\usepackage{amsmath,amssymb,mathtools,unicode-math}
\setmathfont{Latin Modern Math}
\usepackage{enumitem,needspace,xurl,hyperref}
\setlength{\parindent}{0pt}
\setlength{\parskip}{5pt}
\setlength{\emergencystretch}{3em}
\newcommand{\N}{\mathbb{N}}
\newcommand{\Z}{\mathbb{Z}}
\newcommand{\Q}{\mathbb{Q}}
\newcommand{\R}{\mathbb{R}}
\newcommand{\C}{\mathbb{C}}
\newcommand{\F}{\mathbb{F}}
\newcommand{\A}{\mathbb{A}}
\newcommand{\PP}{\mathbb{P}}
\newcommand{\E}{\mathbb{E}}
\newcommand{\eps}{\varepsilon}
\newcommand{\dd}{\,\mathrm{d}}
\begin{document}
\section*{UnsolvedMath ID: 800962}
\textbf{Source problem code:} catalogue entry
\subsection{Definitions and mathematical statement}
Let \(M\) be a smooth compact manifold of dimension \(3\) without boundary, equipped with a reversible Finsler metric \(F\). Thus \(F\) is smooth away from the zero section, positively homogeneous of degree one in each tangent space, and strongly convex, and reversibility means
\[
F(x,v)=F(x,-v)\qquad(x\in M,\ v\in T_xM).
\]
Assume that \(F\) has positive flag curvature: \(K_F(x,y,P)>0\) for every \(x\in M\), every nonzero \(y\in T_xM\), and every two-dimensional plane \(P\subset T_xM\) containing \(y\).

Let \(N\subset M\) be a closed two-dimensional totally geodesic submanifold. Here closed means compact without boundary, and totally geodesic means that geodesics of the induced Finsler metric on \(N\) are also geodesics of \((M,F)\). The draft uses the usual embedded-submanifold convention; the supplied source does not explicitly specify whether immersed submanifolds are also intended.

A closed geodesic is a nonconstant periodic \(F\)-geodesic \(\gamma:\R/T\Z\to M\), with \(T>0\). Is it true that, for every \((M,F)\) and \(N\) satisfying these hypotheses and every such \(\gamma\),
\[
\gamma(\R/T\Z)\cap N\ne\varnothing
\quad\text{or}\quad
\gamma(\R/T\Z)\subseteq N\,?
\]
The second alternative already implies the first under the ordinary set-theoretic meaning of intersection. Both alternatives are retained from the source; no transversality requirement is asserted.
\subsection{Short English statement}
Must every nonconstant closed geodesic on a closed reversible Finsler three-manifold of positive flag curvature meet a given closed totally geodesic surface, including the possibility that it lies entirely in that surface?
\subsection{Sources}
\begin{itemize}
\item Local UnsolvedMath registry, permanent ID \(800962\), source problem code \texttt{catalogue entry}, title: \emph{Open Problems and Questions about Geodesics --- Question 43}. The supplied record attributes the question to Burns and Matveev, with proposed year 2019.
\item Burns, Matveev, \emph{Open Problems and Questions about Geodesics} (2019), Question 43. Supplied source URL: \url{https://arxiv.org/abs/1308.5417}. The record reports access on 2026-10-16 and extraction by \texttt{arxiv\_latex}; these are supplied provenance details, not independently verified access or source verification.
\item The registry labels the problem \texttt{open}. This is the supplied status label, not an independently verified assessment of the literature.
\end{itemize}
\end{document}
